% =====================================================================
%  MODELO EM BRANCO -- uma ação do MOP nos 13 tópicos padrão
%
%  Como usar:
%   1. Copie para capitulos/<acao>/capitulo.tex (ou divida em um arquivo
%      por tópico, como em capitulos/subcomponente-3-2-b/).
%   2. Substitua tudo o que estiver entre << >>.
%
%  Regras:
%   - Texto curto: um parágrafo de 4 a 6 linhas por tópico, que apresenta
%     a figura ou o quadro. O conteúdo fica nos fluxogramas e quadros.
%   - Fluxogramas com \fluxograma{arquivo} (pasta figuras/). Tipos:
%       sequência .......... figuras/fluxograma-captacao-projeto.tex
%       raias .............. figuras/fluxograma-obras.tex
%       linha do tempo ..... figuras/fluxograma-modelo-gestao.tex
%       zigue-zague ........ figuras/fluxograma-abastecimento.tex
%       mapa mental ........ figuras/mapa-mental.tex
%   - Prazos, metas e custos: planilha de custos do PSH-PR.
%   - Tópico 10: descrever as medidas por tema, sem citar as NAS.
%   - Inconsistências: \pendencia{...}.
% =====================================================================
\chapter{<<Nome da ação>>}

\section{Mapa mental da ação}
<<Problema, objetivo, executor e prazo em um parágrafo.>>
% figura: mapa mental (copiar figuras/mapa-mental.tex)

\section{Critérios de seleção}
<<Quem ou onde a ação atende e por quê.>>
% figura: fluxograma de aplicação dos critérios

\section{Mapa das áreas de intervenção}
<<Onde a ação acontece; como o mapa é atualizado.>>
% figura: \includegraphics{figuras/<<mapa>>.pdf}  (ver mapa/gerar_mapa.py)

\section{Responsabilidades das instituições}
<<Papel de cada instituição em um parágrafo.>>
% quadro: roteiro por ação (Ação | Entradas | Responsável | Órgão | Saída)
% quadro: matriz E/A/V

\section{Visão geral da implementação}
<<Blocos da implementação, premissas, financiamento.>>
% figuras: visão geral; ciclo de contratação

\section{<<Etapa estruturante, ex.: Modelo de gestão>>}
\subsection{<<Subetapa, ex.: Trabalho técnico-social>>}

\section{Implementação de <<produto 1>>}
% uma subseção curta por fase, cada uma com seu fluxograma

\section{Implementação de <<produto 2>>}

\section{Conexão com outros subcomponentes do MOP}
% figura: o que a ação recebe e o que entrega

\section{Gestão ambiental e social}
% figura: medidas por fase; quadro: Tema | Como a ação atende | Evidência

\section{Controle de qualidade}
% figura: cadeia de comprovação; quadro: matriz de comprovação por produto;
% quadro de indicadores; fluxo de monitoramento

\section{Matriz de custos da ação}
% quadro: item, executor, fonte, unidade, custo unitário, quantidade,
% total e valores por ano

\section{Cronograma}
% figura: Gantt anual; tabela de metas físicas anuais; remissão ao POA
